\section{Introdução} \label{sec:firstpage}

O ensino e a experimentação em Inteligência Artificial (IA) frequentemente utilizam ambientes simplificados que permitem explorar conceitos fundamentais sem a complexidade inerente aos problemas do mundo real. Nesse contexto, o Mundo do Wumpus destaca-se como um dos exemplos clássicos para o estudo de agentes inteligentes e processos de tomada de decisão, combinando regras relativamente simples com desafios relacionados à percepção parcial do ambiente, raciocínio sob incerteza e planejamento. O conceito tem origem no jogo Hunt the Wumpus, proposto por \cite{yob1975hunt}, com o objetivo de criar um desafio de exploração e sobrevivência em um ambiente desconhecido e dinâmico, envolvendo elementos como armadilhas, um inimigo oculto e a necessidade de localização espacial do jogador. Posteriormente, \cite{russell2021artificial} incorporaram esse cenário ao contexto da Inteligência Artificial, consolidando o Mundo do Wumpus como um ambiente clássico para o estudo de agentes inteligentes, representação do conhecimento e mecanismos de inferência lógica.

No Mundo do Wumpus, o agente navega em um grid parcialmente observável contendo poços, uma criatura hostil e ouro, devendo coletar o ouro e retornar à saída sem ser eliminado. O agente utiliza percepções locais (como brisa para indicar poços adjacentes e fedor par proximidade do Wumpus) para inferir perigos, e dispõe de flechas limitadas para eliminar a criatura, o que exige decisões estratégicas ao longo da exploração.

O desempenho do agente é avaliado por meio de um sistema de pontuação associado às ações executadas durante a simulação, incluindo movimentação, utilização de flechas, coleta do ouro e sobrevivência. Esse mecanismo permite a comparação objetiva entre diferentes estratégias e abordagens de tomada de decisão.

O Mundo do Wumpus é amplamente empregado para demonstrar o comportamento de agentes inteligentes em ambientes parcialmente observáveis e incertos, permitindo a investigação de diferentes estratégias de resolução. Além do estudo de agentes baseados em lógica, o ambiente pode ser adaptado para pesquisas envolvendo ambientes dinâmicos, algoritmos evolutivos, sistemas multiagentes e outras abordagens relacionadas à Inteligência Artificial. Sua flexibilidade faz com que seja utilizado tanto em atividades de ensino quanto em projetos de pesquisa.

Embora o Wumpus World seja amplamente utilizado como benchmark, a ausência de uma plataforma unificada que padronize a configuração de ambientes, a definição de agentes e a coleta de métricas dificulta a comparação objetiva entre diferentes abordagens e a reprodução de experimentos. Pesquisadores frequentemente precisam reimplementar o ambiente ou adaptar código legado, o que consome tempo e introduz vieses.

Diante desse cenário, este trabalho apresenta o projeto “Wumpus Verse”, um protótipo de uma plataforma web que elimina a necessidade de uma instalação local preservando sua flexibilidade e potencial para experimentação. A ferramenta oferece uma interface gráfica amigável e um conjunto de funcionalidades que permitem aos usuários sem experiência em programação criar e personalizar mundos, configurar diferentes tipos de agentes inteligentes, executar experimentos e armazenar os resultados obtidos. Além disso, por funcionar em ambiente web, a plataforma amplia a acessibilidade e fornece um ecossistema integrado para atividades de ensino, aprendizagem e pesquisa envolvendo agentes inteligentes.